\section{When the efficient pattern alternates}

Assume $n=3m$ with $m\ge3$. This case has a different efficient pattern. Here the period is
$p_0=2m$, so every chain in a nontrivial path has length at most
$b=p_0-1$. Lemma~\ref{lem:inherited-chains} forbids four consecutive chains
whose middle two both have length $b$ and whose outer two have
length at least two.

Temporarily suppose that no chain has length one. Away from the ends of
the path, two chains of maximum length cannot then be adjacent. The
largest average allowed by this restriction comes from alternating $b$
and $b-1$. For example, at $n=12$ the candidate sequence is
$7,6,7,6,\ldots$, with mean $13/2$. The old parameter is accordingly
$t_0=2b-1=2p_0-3$.

As before, this efficient repetition returns to an occupied class. The
additional issue is how to count departures from it. A chain of length
one can allow a neighboring pair of maximum lengths that would otherwise
be forbidden. Such a short chain loses many units of length, and some of
that loss must pay for the exception it creates. Counting only the number
of nonalternating positions would discard this information.

\subsection{The return of an alternating pair}

The normalization of one pair is $h_{b-1}h_b$. It has an especially
simple description after changing coordinates: advance seven places
backwards around a circle, then interchange two neighboring symbols.
A transposition either splits one cycle into two or joins two cycles
into one. This gives the order without following a long word of cells.

Define
\begin{equation}\label{eq:R0}
 R_0=\begin{cases}
 2n/7,&7\mid n,\\
 u(n-u),&7\nmid n,
 \end{cases}
\end{equation}
where in the second case $1\le u<n$ and $7u\equiv1\pmod n$.
For $n=12$, the inverse is $u=7$, giving $R_0=35$.

\begin{lemma}\label{lem:alternating-return}
An alternating factor of $2R_0+1$ chains, with lengths $b,b-1$ in
either order, cannot occur in a genuine light segment. Moreover,
$R_0\le n^2/4$.
\end{lemma}

\begin{proof}
Let $s$ be the successor permutation of the circle $(1,3,\ldots,n,2)$,
and put $\tau=(1\ 3)$. The permutations in \eqref{eq:maps} give
$\rho=\tau s$ and $g=\tau s^3$. Since $g$ has order $p_0$,
\[
 M_0=h_{b-1}h_b=\rho g^{-2}\rho g^{-1}
     =\tau s^{-2}\tau s^{-5}\tau.
\]
Conjugation by $s^2\tau$ gives $\tau s^{-7}$, as described above.
If $7\nmid n$, the map $s^{-7}$ is an $n$-cycle. The transposition
cuts it into arcs of lengths $u$ and $n-u$. These lengths are coprime,
so the order is $u(n-u)\le n^2/4$.
If $7\mid n$, the map $s^{-7}$ has seven cycles of length $n/7$.
The transposition joins two and leaves the other five, giving order
$2n/7\le n^2/4$.

Reversing the pair gives a conjugate map with the same order. In an
alternating factor of $2R_0+1$ chains, the first $2R_0$ have outgoing
shortened gaps. Their $R_0$ pairs have combined normalization equal
to the identity, so the following short-run start repeats the first
start class. The post-$B$ forcing rule applies because each such run
has $n-2\ge7$ blocks.
\end{proof}

\subsection{Paying for the breaks in alternation}

We first describe the count. Measure each chain's loss from its maximum
length $b$. The two efficient lengths $b,b-1$ then become the binary
symbols $0,1$, and the efficient pattern becomes an alternating binary
string. We shall remove larger losses, cut at adjacent equal binary
symbols, and apply Lemma~\ref{lem:alternating-return} to the alternating
strings that remain.

In the interior of the path there are two kinds of equal adjacency.
A pair $11$ already contains
extra loss: both chains have fallen below the maximum. A pair $00$ has
no such loss, but the four-chain restriction requires a neighboring
chain of length one. That chain supplies the payment. The following
proof keeps track of these two sources separately.

\begin{proposition}\label{prop:CP0}
For $n=3m$, $m\ge3$, we have
\[
 \CP\left(p_0+\frac12+\frac1{4R_0},\,
 t_0-\frac1{2R_0}\right).
\]
\end{proposition}

\begin{proof}
Consider a path of $q\ge2$ chains. Write $z_i=b-\ell_i$ for the loss
of its $i$th chain, and let $w_i$ indicate whether that chain has
length one. For each four-chain window, at least one of the middle
chains is below maximum length or one of the outer chains has length
one. Thus, for $2\le i\le q-2$,
\[
 f_i:=\ind_{z_i>0}+\ind_{z_{i+1}>0}+w_{i-1}+w_{i+2}\ge1.
\]
Think of each summand as a witness to this local restriction.
An ordinary nonmaximum chain supplies one marker, and a chain of length
one supplies two: it can witness the restriction in a middle position
or in an outer position. Accordingly, put $v_i=\ind_{z_i>0}+w_i$.
Each marker occurs in at most two of the windows counted by $\sum f_i$.
Also $v_i\le z_i$: a chain of length one loses $b-1\ge4$ units,
more than enough to supply its two markers.

There are $q-3$ windows when $q\ge3$. This suggests measuring the
remaining loss by $E=2\sum_i z_i-(q-3)$, after subtracting their
mandatory witnesses. Let $A=\sum_{i=2}^{q-2}(f_i-1)$ count the extra
witnesses. Then
\[
 E=2\sum_i(z_i-v_i)
   +\left(2\sum_i v_i-\sum_{i=2}^{q-2}f_i\right)+A.
\]
All three terms are nonnegative. A nonmaximum chain other than a chain
of length one spends one loss unit on its marker. If its loss is at
least two, at least one unit remains. A chain of length one spends two
of its $b-1$ units, leaving $b-3$. Consequently, if $B_0$ counts the
indices with $z_i\ge2$ and $\ell_i\ne1$, and $W$ counts those with
$\ell_i=1$, then
\begin{equation}\label{eq:defect-payment}
 E\ge2B_0+2(b-3)W+A.
\end{equation}
This is the reserve available to pay for cuts.

Now restrict to the interior indices $2,\ldots,q-1$ and remove those
counted by $B_0$ or $W$, leaving breaks where they stood. The surviving
strings are binary. Cut at every $00$ or $11$ adjacency. Each $11$
contributes at least one extra witness to $A$. Each $00$ needs a
neighboring chain of length one, and a single such chain can explain
at most two $00$ adjacencies. Hence the total number of alternating
strings is at most $B_0+3W+A+1$: a removed chain creates one break,
and a chain of length one may account for two further cuts.

Every remaining string has length at most $2R_0$. Counting the removed
indices as well gives
\begin{equation}\label{eq:alternating-count}
 q-2\le B_0+W+2R_0(B_0+3W+A+1)\le2R_0(E+1).
\end{equation}
For the last inequality, \eqref{eq:defect-payment} pays for each type
of break separately; $b\ge5$ and $R_0\ge1$ suffice for the coefficients.
At $q=3$ the window sums are empty and the argument still applies.
At $q=2$, the identity $E=2\sum_i z_i+1\ge1$ gives
\eqref{eq:alternating-count} directly.

Finally let $\mu=t_0/2=b-1/2$. The total length is
\begin{equation}\label{eq:loss-length}
 \sum_i\ell_i=\mu q+\frac32-\frac E2.
\end{equation}
Inserting \eqref{eq:alternating-count} into \eqref{eq:loss-length}
gives
\[
 \sum_i\ell_i\le
 \left(\mu-\frac1{4R_0}\right)q+2+\frac1{2R_0}
 =p_0+\frac12+\frac1{4R_0}
  +\frac12\left(t_0-\frac1{2R_0}\right)(q-1).
\]
For a single chain, the earlier bound $\ell_1\le p_0$ suffices.
\end{proof}

The average has decreased, but the endpoint allowance has increased
slightly. The count explains both changes: alternating pieces have bounded
length, while the two ends can leave incomplete pieces. Only the first
cost accumulates with the number of chains.
The same proof permits any real upper bound $U\ge R_0$ in place of
$R_0$, because every alternating string has length at most $2R_0\le2U$.
In particular, the simpler choice $U=n^2/4$ is valid.
